\section{Conclusion}

The BAKO Motors Intelligent Supervision System (ISS) project set out to solve five specific operational gaps in the BAKO Motors electric motorcycle prototype: the absence of real-time sensor coverage beyond the battery pack, the lack of a monitoring interface for engineers, the restriction of CAN bus communication to the BMS only, insufficient system documentation for maintenance teams, and the inability to perform remote diagnostics. All five gaps have been addressed by the system delivered in this project.

The hardware layer — an ESP32-based custom two-layer PCB integrating an MCP2515 CAN controller, TJA1051T transceiver, three DS18B20 temperature sensors, two ACS712 current sensors, and a SIM800L GSM/GPRS module — provides a unified data acquisition platform for the entire vehicle. The firmware layer implements a single-threaded cooperative loop that decodes CAN frames from the O'CELL IFS60.8-500-F-E3 BMS, polls external sensors at appropriate rates, and publishes telemetry snapshots over both Wi-Fi and cellular. The cloud layer, built on FastAPI and deployed on a Linux VPS, receives these snapshots via HTTP POST, stores them in SQLite, and streams them to browser clients at up to 10~Hz via WebSocket. The dashboard layer presents all 19 battery cells, four temperature probes, current and voltage measurements, and the vehicle handbrake state in a single real-time web interface with data export capability.

The project was executed across eight two-week Agile sprints, achieving an average velocity of 18 story points per sprint and delivering 144 of the 160 committed story points. The 16 deferred points represent explicitly deprioritised items (OTA firmware update, enclosure design, full-solar field validation) rather than incomplete or broken features. All quantitative V\&V criteria were met: CAN frame reception $> 99\%$ after termination correction, DS18B20 temperature accuracy $\pm 0.5$\textdegree{}C, ACS712 current error $\pm 1.8\%$, voltage measurement error $\pm 0.9\%$, and GSM alert latency $< 15$~s. The dashboard renders at $\geq 5$ frames per second and the cloud pipeline sustained zero snapshot loss over a 1-hour test.

The total project expenditure was approximately 240~TND (60~TND in electronics components and 180~TND in PCB fabrication), representing a 5.9\% underspend against the planned budget, achieved through successful local component sourcing and PCB fabrication negotiation. The board is at revision RECTIFICATION\_2 and is ready for Gerber export and submission to a PCB fabrication house for a production batch.

Seven key findings were identified during the project that will inform future monitoring system designs for Tunisian electric vehicle manufacturers: the criticality of early CAN protocol reverse engineering, the superiority of Wi-Fi AP mode over Station Mode for in-vehicle applications, the accuracy of voltage-derived SOC for LiFePO4 chemistries, the mandatory 2200~\textmu{}F bulk capacitor requirement for the SIM800L GSM/GPRS module, the marginal performance of ACS712-30A at low currents, and the importance of cross-subsystem integration testing from the third sprint onwards.

The BAKO Motors ISS prototype demonstrates that a fully functional, cloud-connected vehicle monitoring system can be designed, fabricated, and validated by an academic engineering team using open-source tools, locally available components, and a total budget under 600~TND. The architecture and codebase are documented, open, and ready to serve as the foundation for the fleet management platform and standalone product roadmap described in Section~\ref{sec:recommendations}.
